% Lösungen Einheit 3 – Klammern auflösen (zusammenbau v1.0)
\kbabschnitt{Klammern auflösen}

\erg{15}{$3a + 17$ (richtig $\rightarrow$ Nr.~17–18 überspringen)}
\erg{16}{a) ein Plus \quad b) ein Minus \quad c) die Zahl $5$ als Faktor \quad d) $-6$ als Faktor}
\erg{17}{a) Klammer weglassen: $5a + (2b - 1) = 5a + 2b - 1$ \quad b) Klammer weglassen: $5a + (2b - 3) = 5a + 2b - 3$ \quad c) Klammer weglassen: $5a + (2b - 4) = 5a + 2b - 4$ \quad d) $5x + 15$ \quad e) Klammer zuerst: $30 - 20 = 10$; Glied für Glied: $30 - 12 - 8 = 10$; Ergebnis: beide Wege ergeben 10. \quad f) $5a - 2b - 4$ \quad g) $7a - 2b + 5 - c$ \quad h) $-2x - 10$ \quad i) $7x + 12$}
\erg{18}{$4x + 23$}
\erg{19}{a) Das Minus vor der Klammer dreht beide Vorzeichen, aus $-5$ wird $+5$. Richtig: Klammer auflösen: $14 - (a - 5) = 14 - a + 5$; zusammenfassen: $= 19 - a$. \quad b) Ja: $5 - (x - 2) = 5 - x + 2 = 7 - x$. \quad c) $5 \cdot (x + 2) = 5x + 10$ \quad d) Term mit Klammer: $9 \cdot (24 + x)$; Klammer auflösen: $9 \cdot (24 + x) = 216 + 9x$. Die Klasse zahlt $216 + 9x$ Euro.}
